\section{How Data Is Found}
\label{sec:routing}

A request for contract $\contract$ originates at some peer and must reach a peer that hosts $\contract$'s state. This section describes the ring topology, the rule that constructs it, the routing logic that uses it, and the bootstrap by which a new peer joins. The section closes with a 24-hour measurement of routing path lengths on the live network.

\subsection{Locations and the Ring}

A peer's location $\ell(p)$ is computed by hashing its external network address into $[0,1)$. A contract's location $\ell(\contract)$ is similarly derived from its content-addressed key. Distance is the wrap-around metric of \cref{sec:primitives}.

This arrangement makes locations stable across restarts that preserve a peer's address. Locations change when a peer is rehomed, when NAT mappings shift, or when a peer is reachable on multiple addresses. An adversary controlling many addresses (a cloud subnet, an IPv6 range, or a small botnet) can grind locations and concentrate them in chosen regions of the ring. Sybil resistance is left to application-level mechanisms such as signed contributions, reputation contracts, and capacity-limited delegation (\cref{sec:trust}).

\subsection{Why Small-World Routing}
\label{sec:smallworld}

The routing strategy derives from Kleinberg's small-world result \citep{kleinberg2000small}. Consider $N$ peers on a ring (or any metric space), each with a constant number of \emph{short-range} links to immediate neighbors and a small number of \emph{long-range} links sampled with probability proportional to $1/d$ in ring distance $d$. Kleinberg proved that a node with knowledge limited to its neighbors can greedily route to any target in $O(\log^2 N)$ hops. The result requires neither global topology knowledge nor flooding, provided the link distribution remains approximately $1/d$.

The $1/d$ distribution is what makes the bound work. Too short and the long-range links don't help cover ground; too long and they overshoot the target. $1/d$ produces, at each scale, a constant probability of a useful long-range hop, so a request roughly halves its remaining distance at every step until the local neighborhood dominates.

Kleinberg's analysis assumes the $1/d$ distribution. Symphony \citep{manku2003symphony} sampled long-range links from a harmonic distribution at join time and retained them. Freenet instead adjusts each peer's neighborhood continuously. A peer directs new \textsf{CONNECT} requests toward under-represented regions of its current connection distribution (\emph{gap-targeting}) and prefers incoming candidates that close large gaps (\emph{Kleinberg-style scoring}). These mechanisms maintain the distribution under churn.

\subsection{Topology Construction}

The \textsf{CONNECT} operation constructs the small-world topology. A joiner chooses a \emph{desired location}, routes a \textsf{CONNECT} message greedily toward it, and receives acceptance decisions along the route and at its terminus. Target selection and acceptance both shape the topology toward Kleinberg's $1/d$ distribution \citep{kleinberg2000small}.

% Ring + gap target. Shows a peer's view of the ring with existing neighbors
% as points on the circle, the largest log-distance gap highlighted as an arc,
% and the gap midpoint marked as the next CONNECT target.
\begin{figure}[t]
\centering
\begin{tikzpicture}[
    peer/.style={circle, draw=black!70, fill=white, inner sep=1.2pt},
    self/.style={circle, draw=black, fill=black, inner sep=1.8pt},
    target/.style={star, star points=5, draw=black!80, fill=yellow!70, inner sep=1.4pt},
    label/.style={font=\scriptsize},
    >=Stealth,
]
  \def\R{2.6}
  % The ring
  \draw[gray!60, thick] (0,0) circle (\R);

  % Self at top (location 0)
  \node[self, label={[label]above:self ($\ell{=}0$)}] (S) at (90:\R) {};

  % Existing neighbors at angles chosen to leave one large gap on the
  % left-hand side of the ring (between 200deg and 310deg, ~ a third of the ring).
  \foreach \a/\name in {62/n1, 38/n2, 12/n3, -10/n4, -40/n5, -75/n6, -130/n7, 165/n8, 130/n9, 110/n10} {
    \node[peer] (\name) at (\a:\R) {};
  }

  % Highlight the large gap between n7 (-130 = 230deg) and n8 (165deg).
  % Span the gap clockwise from n7 to n8 (going through 180deg).
  \draw[red!80, line width=1.8pt] (n7) arc[start angle=-130, end angle=-195, radius=\R];

  % Gap midpoint (-130 + 165 averaged on the short side: midpoint of the arc
  % from -130 going clockwise to -195 is at -162.5 = 197.5deg).
  \node[target, label={[label, text=red!80]left:next \textsf{CONNECT} target}] (T) at (-162.5:\R) {};

  % Bracket-ish annotation pointing to the gap.
  \draw[red!80, ->, thick] ($(-162.5:\R*1.45)$) -- ($(-162.5:\R*1.08)$);
  \node[label, text=red!80, align=center] at (-162.5:\R*1.78) {largest gap\\(log-distance)};

  % Annotation for an existing neighbor distance (just one, to keep clean).
  \draw[blue!50!black, dashed, thin] (S) -- (n3);
  \node[label, text=blue!50!black] at ($(S)!0.5!(n3)$) [above right=-1pt] {$d_3$};

  % Caption-side annotation for the neighbor cloud.
  \node[label, gray!60!black] at (45:\R*1.30) {existing};
  \node[label, gray!60!black] at (45:\R*1.40) {neighbors};

\end{tikzpicture}
\caption{Gap-targeting from a single peer's view. The peer measures
ring distance $d_i$ to each existing neighbor (one such $d_i$ shown
dashed), identifies the largest gap in log-distance space (red arc),
and chooses its midpoint as the next \textsf{CONNECT} destination
(star). The peer thus pulls its own connection distribution toward
$1/d$ as it accumulates neighbors.}
\label{fig:ring-gap}
\end{figure}


\paragraph{Desired location.} A peer with fewer than $K_0$ existing connections targets its own location, with random jitter applied after consecutive acceptance failures so that subsequent attempts probe different ring regions. Joiner locations are derived from network-address prefixes (\cref{sec:trust}), so distinct joiners hit distinct early targets and bootstrap traffic naturally spreads across the gateway's neighborhood. A peer with at least $K_0$ connections instead computes a \emph{gap target}: it measures the ring distance from itself to each existing neighbor, identifies the largest gap in that distribution in log-distance space, and targets the gap's midpoint. Gap-targeting is the active mechanism by which each peer pulls its own connection distribution toward $1/d$ as it accumulates neighbors: the locations under-represented by current connections are exactly the locations the next \textsf{CONNECT} targets. The deployed thresholds are $K_0 = 3$ for the initial-join path and $K_0 = 5$ for the steady-state maintenance loop; the slightly higher maintenance threshold lets the local connection distribution settle before refresh \textsf{CONNECT}s start gap-targeting.

\paragraph{Acceptance.} A peer that receives a \textsf{CONNECT} applies a chain of decisions. If it can forward strictly closer to the target, it forwards; this is the original terminus rule. A peer within a small fraction of the ring of the target (the deployed band is 5\%) may additionally accept with a proximity-scaled probability while still forwarding, giving the joiner more than one nearby connection per request and softening the strict-terminus discipline. When a request reaches its terminus, the receiver's choice depends on its current connection count: well below a small threshold it accepts unconditionally to unblock bootstrap; below the configured minimum-connections target it filters by a Kleinberg-style score on the candidate's location relative to its existing neighbors, with the acceptance floor sliding from generous (early) to selective (near full); at or above the minimum it admits only the best-scoring candidate from a rolling window of recent arrivals. Above the maximum it rejects. When the terminal receiver rejects, the \textsf{CONNECT} may also be re-forwarded \emph{uphill} (away from the target, up to a small budget of additional hops, deployed value 8) so that a joiner whose precise target lies behind a saturated peer can still be picked up by a nearby non-saturated one.

% CONNECT path with accept-and-forward in the proximity band. The joiner
% issues a single CONNECT; the request is greedily forwarded toward the
% joiner's desired location; peers within the proximity band of the target
% may both accept and forward, so the joiner ends up with more than one
% neighbor connection from a single CONNECT.
\begin{figure}[t]
\centering
\begin{tikzpicture}[
    peer/.style={circle, draw=black!70, fill=white, inner sep=2.2pt, font=\scriptsize},
    joiner/.style={circle, draw=black, fill=blue!18, inner sep=2.5pt, font=\scriptsize},
    terminus/.style={circle, draw=black, fill=green!28, inner sep=2.5pt, font=\scriptsize},
    band/.style={circle, draw=black, fill=orange!30, inner sep=2.5pt, font=\scriptsize},
    fwd/.style={->, >=Stealth, gray!70!black, thick},
    accept/.style={->, >=Stealth, blue!60!black, thick, dashed},
    label/.style={font=\scriptsize},
    >=Stealth,
]
  % Layout: joiner on left, then 5 forwarding hops, last 2 inside proximity band of target.
  \node[joiner]    (J)  at (0,0)    {J};
  \node[peer]      (P1) at (1.6,0)  {};
  \node[peer]      (P2) at (3.2,0)  {};
  \node[peer]      (P3) at (4.8,0)  {};
  \node[band]      (P4) at (6.4,0)  {};
  \node[band]      (P5) at (8.0,0)  {};
  \node[terminus]  (T)  at (9.6,0)  {T};

  % Shaded band behind P4..T to make the proximity band visually obvious.
  \begin{scope}[on background layer]
    \fill[orange!12, rounded corners=4pt] ($(P4.south west)+(-0.25,-0.45)$)
                                          rectangle
                                          ($(T.north east)+(0.25,0.45)$);
  \end{scope}

  % Greedy forwarding chain
  \draw[fwd] (J)  -- (P1) node[midway, above, label] {fwd};
  \draw[fwd] (P1) -- (P2);
  \draw[fwd] (P2) -- (P3);
  \draw[fwd] (P3) -- (P4);
  \draw[fwd] (P4) -- (P5);
  \draw[fwd] (P5) -- (T)  node[midway, above, label] {fwd};

  % Accept edges from band peers + terminus back to joiner. Vary the bend
  % so they don't pile on top of each other.
  \draw[accept] (P4) to[bend right=55] (J);
  \draw[accept] (P5) to[bend right=45] (J);
  \draw[accept] (T)  to[bend right=38] (J);

  % Labels
  \node[label] at ($(J)+(0,-0.42)$) {joiner};
  \node[label] at ($(T)+(0,-0.42)$) {target $t$};
  \node[label, gray!55!black] at ($(P2)!0.5!(P3)+(0,0.4)$) {greedy forwarding};

  % Proximity band label, below the shaded region.
  \node[label, text=orange!80!black, align=center]
       at ($(P4)!0.5!(T)+(0,-2.45)$)
       {proximity band of $t$ ($\sim 5\%$ of ring):\\
        these peers \emph{accept and forward}};

  % "accept (terminus)" label inline with the rightmost bend.
  \node[label, text=green!42!black, align=center]
       at ($(T)+(-0.5,-1.65)$) {terminus\\accept};

\end{tikzpicture}
\caption{A single \textsf{CONNECT} from joiner $J$ to its desired
location $t$. The request is greedily forwarded toward $t$ (solid
arrows). Once the route enters the proximity band of $t$ (the
deployed band is $\sim$5\% of the ring, shown shaded), each receiver
may \emph{additionally} accept the joiner with a proximity-scaled
probability while still forwarding the request onward (dashed return
edges). The terminus also applies its acceptance chain
(\S\ref{sec:routing}). Net result: one \textsf{CONNECT} request can
establish multiple connections close to $t$, not just one at the
terminus.}
\label{fig:connect-path}
\end{figure}


\paragraph{Formation of long-range edges.} Gap-targeting and Kleinberg-score acceptance jointly produce the long-range component of the connection distribution. A peer with mostly local neighbors has its largest log-distance gap at long range, so its next \textsf{CONNECT} targets a distant terminus. The receiver favors candidates whose locations close large gaps in its own distribution.

\paragraph{Bounded fan-out.} Admission control bounds the connection count between configurable minimum and maximum (deployed defaults 25 and 200).

\subsection{Greedy Forwarding}

A request from peer $p$ to target location $t$ proceeds as follows. Let $N(p)$ be $p$'s current set of neighbors. If some neighbor $q \in N(p)$ is strictly closer to $t$ on the ring than $p$ itself, forward the request to the closest such $q$; otherwise $p$ is a terminus for the request and serves it locally.

Each request carries a \emph{hops-to-live} (HTL) counter, deployed maximum 10, that decrements on each forward and bounds how far a request travels: when HTL reaches zero the current peer stops forwarding and answers from local state or reports a miss. HTL is a termination bound only. Every forward is greedy --- performance-ranked once the adaptive layer of the next subsection has enough history --- and a peer that is near $t$ but has not cached the contract forwards to the next-best candidate under a skip-list that prevents revisiting, rather than switching to a random walk. (Topology-mixing randomness lives in the \textsf{CONNECT} path of \cref{sec:routing}, not in request forwarding.)

\subsection{Adaptive Routing}

Ring distance alone does not identify the fastest path because neighbors may differ by an order of magnitude or more in bandwidth, RTT, and uptime. The router therefore selects the next hop by predicted total cost.

For each neighbor $q$, the local peer maintains a history of routing events of the form $(d, \tau, b, f)$, where $d$ is the ring distance from $q$ to the request's target, $\tau$ is the observed time-to-first-byte, $b$ is the observed transfer rate, and $f$ is a failure indicator. From this history three estimators are fit:
\[
  \widehat{\tau}_q(d), \quad \widehat{p}_q(d), \quad \widehat{b}_q(d),
\]
the predicted latency, failure probability, and throughput as functions of distance-to-target. Each is fit by \emph{isotonic regression} \citep{barlow1972statistical, deleeuw2009isotone}, the standard non-parametric estimator for a monotone function.

The choice of isotonic regression encodes a structural prior we are willing to commit to (latency and failure are non-decreasing in distance-to-target, on average) without committing to a parametric family for the shape of the curve. The pool-adjacent-violators algorithm fits each curve in time logarithmic in history size; old events age out so the model tracks current conditions.

Because the estimators are indexed by distance to the target, each fitted curve combines the neighbor's link quality (RTT, bandwidth, and reliability) with the quality of onward paths through that neighbor. Onward-path quality changes more rapidly under churn because it depends on peers beyond the immediate link. An aging window removes stale observations as new routing events re-score the neighbor.

Given a target $t$ for a new request, the router selects the neighbor that minimizes an expected-cost score combining the three estimators. A neighbor with too little history falls back to greedy-by-distance until enough events have accumulated. Predictions between fitted points use linear interpolation, and beyond the training-data domain the model extrapolates linearly.

\paragraph{Scope of adaptive routing.} Once sufficient history is available, ordinary requests (\textsf{GET}, \textsf{PUT}, \textsf{SUBSCRIBE}, and \textsf{UPDATE}) select neighbors by predicted cost. \textsf{CONNECT} retains a hard greedy filter: only neighbors strictly closer to the target than the current peer are eligible. A separate isotonic estimator trained on accept/reject outcomes ranks those candidates. The filter keeps topology-construction traffic moving toward its target; ordinary requests can use the broader candidate set.

% Adaptive routing: per-neighbor isotonic-regression models of latency vs
% distance-to-target. The router picks the next hop by combining the
% predicted latency, failure probability, and throughput from each
% neighbor's model -- which can differ from greedy-by-distance because a
% near-but-slow neighbor can lose to a farther-but-fast one.
\begin{figure}[t]
\centering
\begin{tikzpicture}[
    axis/.style={->, >=Stealth, thick, black!75},
    curveA/.style={blue!65!black, thick},
    curveB/.style={orange!75!black, thick},
    curveC/.style={green!55!black, thick},
    dot/.style={circle, fill, inner sep=1pt},
    label/.style={font=\scriptsize},
    >=Stealth,
]
  % Axes (single panel: latency vs distance).
  \draw[axis] (0,0) -- (7.0,0) node[right, label] {distance to target $d$};
  \draw[axis] (0,0) -- (0,3.6) node[above, label, align=center]
       {predicted\\latency $\widehat{\tau}_q(d)$};

  % Three neighbor curves. Each is a piecewise-constant step function (the
  % shape isotonic regression produces). Non-decreasing in d. Lower curve =
  % faster neighbor.

  % Neighbor q_A: a fast/long-haul neighbor -- low latency even at large d.
  \draw[curveA] (0.0,0.55) -- (1.2,0.55) -- (1.2,0.75) -- (2.6,0.75)
                -- (2.6,0.95) -- (4.0,0.95) -- (4.0,1.15) -- (5.4,1.15)
                -- (5.4,1.50) -- (6.7,1.50);
  \node[label, text=blue!65!black, anchor=west] at (6.75,1.55) {$\widehat{\tau}_{q_A}$ (fast)};

  % Neighbor q_B: average -- moderate latency rising more steeply.
  \draw[curveB] (0.0,1.10) -- (0.9,1.10) -- (0.9,1.45) -- (2.0,1.45)
                -- (2.0,1.85) -- (3.3,1.85) -- (3.3,2.25) -- (4.8,2.25)
                -- (4.8,2.70) -- (6.7,2.70);
  \node[label, text=orange!75!black, anchor=west] at (6.75,2.75) {$\widehat{\tau}_{q_B}$ (average)};

  % Neighbor q_C: high-latency, even close.
  \draw[curveC] (0.0,1.50) -- (0.8,1.50) -- (0.8,1.85) -- (1.6,1.85)
                -- (1.6,2.30) -- (2.6,2.30) -- (2.6,2.75) -- (3.8,2.75)
                -- (3.8,3.20) -- (5.2,3.20) -- (5.2,3.40) -- (6.7,3.40);
  \node[label, text=green!55!black, anchor=west] at (6.75,3.20) {$\widehat{\tau}_{q_C}$ (slow)};

  % Mark a particular distance d_t (the distance to a specific target t).
  \draw[gray!60, thin, dashed] (3.6,0) -- (3.6,3.4)
       node[label, gray!45!black, anchor=south] at (3.6,3.4) {target distance $d_t$};
  % Dots where the curves meet d_t.
  \node[dot, blue!65!black] at (3.6,0.95)  {};
  \node[dot, orange!75!black] at (3.6,2.25) {};
  \node[dot, green!55!black]  at (3.6,2.75) {};

  % Annotation: q_A is chosen.
  \node[label, gray!35!black, align=center]
       at (5.6,0.40) (choosenote) {choose $q_A$ for\\target at $d_t$};
  \draw[->, >=Stealth, gray!50!black, thin]
       (choosenote.west) -- (3.75,0.95);

\end{tikzpicture}
\caption{Adaptive routing. Each neighbor $q$ carries a non-parametric,
non-decreasing model $\widehat{\tau}_q(d)$ of predicted latency as a
function of ring distance to the target, fit by isotonic regression
on the history of routing events. (Two more such models per neighbor
track predicted failure probability and throughput; one is shown for
clarity.) For a request to a target at distance $d_t$, the router
reads each neighbor's prediction at $d_t$ (dots) and picks the
neighbor that minimizes the combined expected cost. The chosen
neighbor need not be the one topologically closest to the target.}
\label{fig:adaptive-routing}
\end{figure}


\subsection{Lifecycles}

The five operations of \cref{sec:primitives} are realized as transaction state machines coordinated by the local core. Each is assigned a transaction id, tracked through completion, and supports timeout, retry, and composite parent/child relationships.

\paragraph{\textsf{CONNECT}.} As above. A successful \textsf{CONNECT} establishes a single neighbor relationship. A joining peer issues an initial sequence of \textsf{CONNECT} operations to populate its first $K_0$ neighbors using own-location targeting; once it has enough connections for gap analysis, all subsequent \textsf{CONNECT} operations use gap-target locations. Topology growth and refresh are driven by additional \textsf{CONNECT} operations issued from a maintenance loop, not as a side effect of ordinary request traffic.

\paragraph{\textsf{PUT}.} A \textsf{PUT} carries a contract $\contract$ and an initial state $\sigma_0$. The originating peer routes the request toward $\ell(\contract)$. A peer on the routed path checks $\valid_\contract(\sigma_0)$ and hosts the contract; routing and hosting co-occur, and there is no separate cache tier beneath hosting. A peer whose storage budget is already exhausted may refuse the write. Peers in the contract's neighborhood receive the state along the same path and replicate it, and a peer that begins hosting announces the fact to its neighbors (\cref{sec:updates}).

\paragraph{\textsf{GET}.} A \textsf{GET} carries a key $k$. The originating peer routes toward the corresponding location. Any peer along the path that currently hosts the contract may answer with its current state $\sigma$ and parameters; the response travels back along the request path, optionally caching at intermediate peers.

A route that terminates without finding a host does not immediately report a miss. The terminus consults its neighbors' host advertisements (\cref{sec:updates}) and forwards to a neighbor that holds the contract, if one exists. Greedy forwarding reaches the relevant ring neighborhood; the advertisement mesh supplies the final step when the selected peer is not a host. \textsf{SUBSCRIBE} terminates the same way.

The requester can verify on arrival that the served code and parameters hash to $k$. \emph{State} is not directly verifiable from the key: the key is derived from $\mathrm{code}$ and $\mathrm{params}$, not from state, and state is mutable. This is a meaningful difference from immutable content-addressed stores such as IPFS, where $\mathsf{get}(k)$ returns a value whose own hash equals $k$. Readers carrying that mental model should note that the guarantee does not transfer to Freenet contracts. State is trusted only to the extent that the contract's validity predicate accepts it; replication across several peers near $\ell(\contract)$ is the structural defense against a single peer serving a stale or doctored state.

\paragraph{\textsf{UPDATE}.} An \textsf{UPDATE} carries a key $k$ and a candidate state $u \in \statespace_\contract$. It is routed toward $\ell(\contract)$. Each peer along the path that already hosts the contract merges $u$ into its local copy subject to the validity check, then propagates the change outward through the subscription tree (\cref{sec:updates}).

\paragraph{\textsf{SUBSCRIBE}.} A \textsf{SUBSCRIBE} establishes a lease on future updates. The subscriber routes the request toward $\ell(\contract)$ and the first hosting peer responds; intermediate peers along the path record the subscription. Subscriptions are leased and renewed periodically (\cref{sec:updates}).

\subsection{A 24-Hour Snapshot of Routing on the Deployed Network}
\label{sec:routing-measurement}

% Live connection-distance distribution, snapshotted from the deployed
% network. n = 2274 connections across 396 active peers (snapshot
% 2026-05-25, source: the OTel-backed dashboard at nova.locut.us:3133).
%
% For each connection, ring distance = min(|loc_a - loc_b|, 1 - |loc_a - loc_b|).
% Linear histogram, 25 bins of width 0.02. The dashed curve is a
% Kleinberg 1/d reference with the same total mass.
%
% Coordinates are precomputed in cm (x: distance*22, y: count*5/450).
\begin{figure}[t]
\centering
\begin{tikzpicture}[
    bar/.style={fill=blue!22, draw=blue!55!black, line width=0.3pt},
    refcurve/.style={red!70!black, thick, dashed},
    axis/.style={->, >=Stealth, thick, black!75},
    tick/.style={black!75},
    label/.style={font=\scriptsize},
    gridline/.style={gray!20, very thin},
]
  % Background y-grid
    \draw[gridline] (0,1.1111) -- (11.0000,1.1111);
    \draw[gridline] (0,2.2222) -- (11.0000,2.2222);
    \draw[gridline] (0,3.3333) -- (11.0000,3.3333);
    \draw[gridline] (0,4.4444) -- (11.0000,4.4444);

  % Axes
  \draw[axis] (0,0) -- (11.4,0) node[right, label] {connection distance $d$};
  \draw[axis] (0,0) -- (0,5.4) node[above, label] {connections};

  % x-ticks
    \draw[tick] (0.0000,0) -- (0.0000,-0.1) node[below, label, yshift=-1pt] {$0.0$};
    \draw[tick] (2.2000,0) -- (2.2000,-0.1) node[below, label, yshift=-1pt] {$0.1$};
    \draw[tick] (4.4000,0) -- (4.4000,-0.1) node[below, label, yshift=-1pt] {$0.2$};
    \draw[tick] (6.6000,0) -- (6.6000,-0.1) node[below, label, yshift=-1pt] {$0.3$};
    \draw[tick] (8.8000,0) -- (8.8000,-0.1) node[below, label, yshift=-1pt] {$0.4$};
    \draw[tick] (11.0000,0) -- (11.0000,-0.1) node[below, label, yshift=-1pt] {$0.5$};

  % y-ticks
    \draw[tick] (0,0.0000) -- (-0.1,0.0000) node[left, label, xshift=-1pt] {0};
    \draw[tick] (0,1.1111) -- (-0.1,1.1111) node[left, label, xshift=-1pt] {100};
    \draw[tick] (0,2.2222) -- (-0.1,2.2222) node[left, label, xshift=-1pt] {200};
    \draw[tick] (0,3.3333) -- (-0.1,3.3333) node[left, label, xshift=-1pt] {300};
    \draw[tick] (0,4.4444) -- (-0.1,4.4444) node[left, label, xshift=-1pt] {400};

  % Bars (linear histogram, 25 bins of width 0.02)
    \draw[bar] (0.0000,0) rectangle (0.4400,4.5667);
    \draw[bar] (0.4400,0) rectangle (0.8800,3.2222);
    \draw[bar] (0.8800,0) rectangle (1.3200,2.7556);
    \draw[bar] (1.3200,0) rectangle (1.7600,1.9222);
    \draw[bar] (1.7600,0) rectangle (2.2000,1.5333);
    \draw[bar] (2.2000,0) rectangle (2.6400,1.1889);
    \draw[bar] (2.6400,0) rectangle (3.0800,1.0778);
    \draw[bar] (3.0800,0) rectangle (3.5200,1.0000);
    \draw[bar] (3.5200,0) rectangle (3.9600,0.9333);
    \draw[bar] (3.9600,0) rectangle (4.4000,0.9000);
    \draw[bar] (4.4000,0) rectangle (4.8400,0.6000);
    \draw[bar] (4.8400,0) rectangle (5.2800,0.6556);
    \draw[bar] (5.2800,0) rectangle (5.7200,0.5111);
    \draw[bar] (5.7200,0) rectangle (6.1600,0.4444);
    \draw[bar] (6.1600,0) rectangle (6.6000,0.3889);
    \draw[bar] (6.6000,0) rectangle (7.0400,0.3556);
    \draw[bar] (7.0400,0) rectangle (7.4800,0.4444);
    \draw[bar] (7.4800,0) rectangle (7.9200,0.4667);
    \draw[bar] (7.9200,0) rectangle (8.3600,0.3556);
    \draw[bar] (8.3600,0) rectangle (8.8000,0.5222);
    \draw[bar] (8.8000,0) rectangle (9.2400,0.3556);
    \draw[bar] (9.2400,0) rectangle (9.6800,0.4333);
    \draw[bar] (9.6800,0) rectangle (10.1200,0.2667);
    \draw[bar] (10.1200,0) rectangle (10.5600,0.1444);
    \draw[bar] (10.5600,0) rectangle (11.0000,0.2222);

  % 1/d reference (same total mass, bin midpoints, d_min clamp = 0.005).
  % Clip to the plot rectangle so the asymptote near d=0 doesn't escape.
  \begin{scope}
    \clip (0,0) rectangle (11.0,5.0);
    \draw[refcurve] plot coordinates {
      (0.2200,7.6056)
      (0.6600,3.8028)
      (1.1000,2.2244)
      (1.5400,1.5789)
      (1.9800,1.2244)
      (2.4200,1.0000)
      (2.8600,0.8456)
      (3.3000,0.7322)
      (3.7400,0.6467)
      (4.1800,0.5778)
      (4.6200,0.5233)
      (5.0600,0.4778)
      (5.5000,0.4389)
      (5.9400,0.4067)
      (6.3800,0.3789)
      (6.8200,0.3544)
      (7.2600,0.3322)
      (7.7000,0.3133)
      (8.1400,0.2967)
      (8.5800,0.2811)
      (9.0200,0.2678)
      (9.4600,0.2556)
      (9.9000,0.2433)
      (10.3400,0.2333)
      (10.7800,0.2244)
    };
  \end{scope}

  % Median guide line at d = 0.0819
  \draw[gray!50!black, thin, dashed] (1.8018,0) -- (1.8018,4.7)
       node[label, anchor=south west, gray!40!black] at (1.8018,4.7) {median $d = 0.082$};

  % Legend
  \node[label, anchor=north east, fill=white, draw=black!50,
        rounded corners=2pt, inner sep=4pt, align=left]
       at (10.9,4.9) {
       \textcolor{blue!55!black}{$\blacksquare$}\ empirical (live network)\\
       \textcolor{red!70!black}{- -}\ $1/d$ reference (Kleinberg)};

\end{tikzpicture}
\caption{Distribution of per-connection ring distances on the deployed
network. Each of the $n = 2274$ live connections across $N = 396$
active peers (snapshot taken 2026-05-25, source: the telemetry
collector at \texttt{nova.locut.us:3133}) contributes one data point;
bins are linear, width $0.02$. The dashed curve is a Kleinberg $1/d$
reference normalized to the same total mass. The empirical
distribution is monotonically decreasing in $d$ at short range and
stays above the $1/d$ reference in the long-range tail, consistent
with the active gap-targeting and Kleinberg-style acceptance of
\cref{sec:routing} steering each peer's neighborhood toward the
$1/d$ shape Kleinberg's analysis requires. Median connection
distance is $0.082$; the maximum half-ring distance is $0.5$.}
\label{fig:connection-distance}
\end{figure}


This subsection records a single observation of routing behavior on the deployed network. It is a sanity check, not a scaling study. We also include a snapshot of the connection-distance distribution (\cref{fig:connection-distance}), a direct empirical check on the small-world topology described above. The 24-hour window ends 2026-05-24. Network size in that window was 443 distinct active peers in the last 10 minutes, with 604 distinct peer identifiers seen across the full 24 hours (the difference attributable to churn). All peers run with telemetry enabled in the current deployment.

For each \textsf{GET} and \textsf{PUT} transaction we compute a \emph{path length}: the number of distinct peers seen emitting the corresponding \textsf{*\_request} event for that transaction id, which is the number of peers visited (hops $+ 1$). Statistics over operations whose telemetry reached completion in the window:\footnote{\textsf{UPDATE} (median path length 2, $n = 386$) and \textsf{SUBSCRIBE} (median 1, $n = 5{,}359$) propagate primarily through pre-existing subscription trees rather than fresh greedy routes, so their path lengths reflect tree depth from the subscriber to the nearest hosting peer, not routing-layer hop counts. They are excluded from the table here.}

\begin{center}
\small
\begin{tabular}{lrrrrr}
\toprule
Operation & $n$ & mean & median & p95 & p99 \\
\midrule
\textsf{GET} & 1{,}302 & 6.84 & 7 & 11* & 11* \\
\textsf{PUT} &    158 & 6.01 & 7 & 10* & 10* \\
\bottomrule
\end{tabular}
\end{center}

\noindent\textit{*The deployed hops-to-live cap is 10, so any route that would otherwise visit 12 or more peers is dropped before it can appear in the table. The reported p95 and p99 for \textsf{GET} are therefore at the cap itself; they record where the tail is truncated, not where it actually lives. We are not directly observing the upper tail of the path-length distribution.}

The measurement has two further limitations. First, it covers one observation at one network size. At $N \approx 443$, $\log_2 N \approx 8.8$; the observed median path length of 7 peers (6 hops) for \textsf{GET} is consistent with logarithmic-diameter routing, but a single data point at a single $N$ cannot distinguish $O(\log n)$ from $O(\sqrt{n})$, $O(n^{0.4})$, or a constant. Second, the table does not compare against any baseline, so it does not isolate the contribution of adaptive routing from greedy-by-distance.

A controlled simulation sweep across orders of magnitude in $N$ is needed to test the routing claims and compare greedy-by-distance with adaptive routing. The deterministic harness of \cref{sec:status} currently validates correctness and convergence at the tens-of-nodes scale, but it does not emit the required hop-count-versus-$N$ curve. Extending it to do so is listed as an open problem in \cref{sec:status}.
